\documentclass[10pt,a4paper]{article}
\usepackage[utf8]{inputenc}
\usepackage[T1]{fontenc}
\usepackage[english]{babel}
\usepackage{geometry}
\geometry{margin=1.8cm}
\usepackage{enumitem}
\usepackage{booktabs}
\usepackage{fancyhdr}
\usepackage{lastpage}
\usepackage{tcolorbox}
\usepackage{hyperref}

\pagestyle{fancy}
\fancyhf{}
\rfoot{\thepage/\pageref{LastPage}}
\lhead{\small TOEFL ITP 2026 --- Day 5}
\rhead{\small MPrometeo}
\renewcommand{\headrulewidth}{0.4pt}
\setlength{\parindent}{0pt}
\setlength{\parskip}{5pt}

\begin{document}

\begin{center}
{\LARGE \textbf{DAY 5 --- Tuesday Mar 3}}\\[4pt]
{\large Listening Intensive}\\[2pt]
{\normalsize Theme: Aldana on AI + Elementary OS vs Linux Distros}\\[2pt]
\rule{\textwidth}{0.8pt}
\end{center}

% ============================================================
\section*{Grammar: Conjunctions \& Connectors}
% ============================================================

\subsection*{Contrast Connectors}

\begin{center}
\begin{tabular}{lll}
\toprule
\textbf{Connector} & \textbf{Structure} & \textbf{Example} \\
\midrule
although/though & + clause & Although Linux is free, many prefer macOS. \\
despite/in spite of & + noun/gerund & Despite being free, Linux struggles. \\
however & sentence. However, sentence. & It's free. However, it lacks apps. \\
nevertheless & sentence. Nevertheless, ... & It crashed. Nevertheless, I use it. \\
whereas & + clause (comparison) & macOS is polished, whereas Linux is flexible. \\
while & + clause & While elementary OS is elegant, it's niche. \\
\bottomrule
\end{tabular}
\end{center}

\textbf{TOEFL Trap:} ``Despite'' NEVER takes a clause directly.
\begin{itemize}[nosep]
\item WRONG: ``Despite it is free, nobody uses it.''
\item RIGHT: ``Despite being free, nobody uses it.''
\item RIGHT: ``Although it is free, nobody uses it.''
\end{itemize}

\subsection*{Cause/Effect Connectors}

\begin{center}
\begin{tabular}{lll}
\toprule
\textbf{Connector} & \textbf{Structure} & \textbf{Example} \\
\midrule
because/since/as & + clause & Because AI evolves fast, jobs change. \\
because of/due to & + noun & Due to rapid AI growth, jobs are changing. \\
therefore/thus & sentence. Therefore, ... & AI is powerful. Therefore, we must regulate it. \\
consequently & sentence. Consequently, ... & Data leaked. Consequently, trust dropped. \\
so...that & so + adj + that + clause & AI is so fast that humans can't compete. \\
\bottomrule
\end{tabular}
\end{center}

\subsection*{Addition Connectors}

\begin{itemize}[nosep]
\item \textbf{furthermore / moreover / in addition} = formal ``also''
\item \textbf{not only...but also} = emphasis (triggers inversion if first)
\item \textbf{as well as} = does NOT change subject-verb agreement
\end{itemize}

% ============================================================
\section*{Reading Passage 1: Aldana on AI ($\sim$220 words)}
% ============================================================

\begin{tcolorbox}[colback=white, colframe=black,
    title={Based on Aldana: ``Is AI Frying Our Brains?'' (Oct 3, 2025)}]
\small
In his October 2025 column, Maximino Aldana posed a question that
many scientists have been asking privately: is artificial intelligence
making us collectively dumber? His argument was not the typical
technophobic complaint. As a complexity scientist, Aldana approached
the issue from an information-processing perspective.

The human brain, Aldana noted, evolved to solve problems through
effort --- through the friction of thinking. When we outsource that
effort to AI systems that generate instant answers, we are not
merely saving time; we are atrophying the very cognitive muscles
that make us capable of critical thought. It is analogous to using
a wheelchair when your legs work perfectly: convenient in the short
term, debilitating in the long term.

However, Aldana was careful to distinguish between different uses
of AI. Using AI as a tool to amplify human capability --- the way a
telescope amplifies vision --- is fundamentally different from using
it as a replacement for thinking. The danger lies not in the
technology itself but in the laziness it enables.

His conclusion was characteristically sharp: ``The question is not
whether AI is intelligent. The question is whether we will remain
intelligent enough to use it wisely.'' For a scientist who studies
complex systems, the irony is that AI --- our most complex creation
--- may be simplifying us.
\end{tcolorbox}

\subsection*{Questions}

\begin{enumerate}[label=\textbf{\arabic*.}]
\item What is Aldana's main argument?
\begin{enumerate}[label=(\Alph*)]
\item AI should be banned completely
\item AI may weaken our critical thinking if used as a replacement for effort
\item AI is always beneficial
\item Scientists should not use AI
\end{enumerate}

\item The wheelchair analogy illustrates:
\begin{enumerate}[label=(\Alph*)]
\item That AI helps disabled people
\item The danger of convenience replacing necessary effort
\item That wheelchairs are better than walking
\item Nothing relevant to AI
\end{enumerate}

\item ``Atrophying'' is closest in meaning to:
\begin{enumerate}[label=(\Alph*)]
\item strengthening \item weakening from disuse
\item expanding \item replacing
\end{enumerate}

\item According to Aldana, the real danger of AI is:
\begin{enumerate}[label=(\Alph*)]
\item The technology itself
\item The laziness it enables in humans
\item Its cost
\item Its speed
\end{enumerate}

\item The word ``characteristically'' implies:
\begin{enumerate}[label=(\Alph*)]
\item surprisingly \item in a way typical of him
\item accidentally \item quietly
\end{enumerate}
\end{enumerate}

\textbf{Answers:} 1-B, 2-B, 3-B, 4-B, 5-B

% ============================================================
\section*{Reading Passage 2: Elementary OS ($\sim$200 words)}
% ============================================================

\begin{tcolorbox}[colback=white, colframe=black]
\small
In the crowded landscape of Linux distributions, elementary OS
occupies a peculiar position. Unlike Ubuntu, which prioritizes
broad compatibility, or Arch Linux, which appeals to users who enjoy
configuring every detail, elementary OS is designed around a single
principle: simplicity through design.

Its interface, called Pantheon, is deliberately reminiscent of
macOS --- clean, minimal, and opinionated. The developers make
choices \emph{for} the user rather than offering endless options.
This philosophy, borrowed from Apple, has earned elementary OS both
passionate fans and fierce critics. Supporters argue that it
demonstrates how Linux can be beautiful and accessible to non-technical
users. Detractors claim it sacrifices the freedom that makes Linux
valuable in the first place.

The distribution also pioneered a ``pay what you want'' model for
its downloads and app store, challenging the assumption that open-source
software must always be completely free. This decision was
controversial within the Linux community, where free-as-in-freedom
is often conflated with free-as-in-price.

Despite its small market share, elementary OS has influenced the
broader Linux ecosystem. Other distributions have adopted similar
design principles, and the conversation about user experience in
open-source software has shifted considerably since elementary OS
first appeared in 2011.
\end{tcolorbox}

\subsection*{Questions}

\begin{enumerate}[label=\textbf{\arabic*.}]
\item Elementary OS primarily differs from other Linux distros in its:
\begin{enumerate}[label=(\Alph*)]
\item Programming language \item Focus on design simplicity
\item Server capabilities \item Gaming performance
\end{enumerate}

\item ``Opinionated'' in context means:
\begin{enumerate}[label=(\Alph*)]
\item Rude \item Making design decisions for the user
\item Open to customization \item Political
\end{enumerate}

\item ``Conflated'' most likely means:
\begin{enumerate}[label=(\Alph*)]
\item separated \item confused or merged together
\item celebrated \item ignored
\end{enumerate}

\item What was controversial about elementary OS?
\begin{enumerate}[label=(\Alph*)]
\item Its use of Linux
\item Its ``pay what you want'' download model
\item Its color scheme
\item Its name
\end{enumerate}

\item It can be inferred that the author:
\begin{enumerate}[label=(\Alph*)]
\item Dislikes elementary OS
\item Views it as small but influential in the Linux world
\item Thinks it will replace macOS
\item Believes Linux should never have design focus
\end{enumerate}
\end{enumerate}

\textbf{Answers:} 1-B, 2-B, 3-B, 4-B, 5-B

% ============================================================
\section*{Connector Exercises (10 questions)}
% ============================================================

\begin{enumerate}[label=\textbf{\arabic*.}, leftmargin=*]

\item \rule{2cm}{0.4pt} elementary OS resembles macOS visually, it
runs on a completely different kernel.
\begin{enumerate}[label=(\Alph*)]
\item Despite \item Although \item Because \item Due to
\end{enumerate}

\item Linux is powerful; \rule{2cm}{0.4pt}, many casual users find
it intimidating.
\begin{enumerate}[label=(\Alph*)]
\item therefore \item moreover \item however \item because
\end{enumerate}

\item \rule{2cm}{0.4pt} its small market share, elementary OS has
influenced other distributions.
\begin{enumerate}[label=(\Alph*)]
\item Although \item Despite \item However \item While is
\end{enumerate}

\item AI is \rule{2cm}{0.4pt} powerful \rule{2cm}{0.4pt} it can
generate human-like text in seconds.
\begin{enumerate}[label=(\Alph*)]
\item such...that \item so...that \item too...that \item enough...that
\end{enumerate}

\item Aldana supports AI as a tool; \rule{2cm}{0.4pt}, he warns
against using it as a thinking replacement.
\begin{enumerate}[label=(\Alph*)]
\item furthermore \item nevertheless \item because \item since
\end{enumerate}

\item \textit{Find the error:}
(A)\textbf{Despite} AI (B)\textbf{is} a powerful tool, many
(C)\textbf{scientists} remain (D)\textbf{cautious} about its impact.

\item \textit{Find the error:}
The software is free; (A)\textbf{moreover}, (B)\textbf{it's}
interface is (C)\textbf{cleaner} than (D)\textbf{most} alternatives.

\item Not only \rule{2cm}{0.4pt} elementary OS beautiful, but it is
also remarkably stable.
\begin{enumerate}[label=(\Alph*)]
\item is \item it is \item does \item has
\end{enumerate}

\item \rule{2cm}{0.4pt} AI evolves rapidly, regulations struggle to
keep pace.
\begin{enumerate}[label=(\Alph*)]
\item Because \item Despite \item Although being \item Due to
\end{enumerate}

\item The Linux community values freedom; \rule{2cm}{0.4pt}, some
users prefer the simplicity of curated systems like elementary OS.
\begin{enumerate}[label=(\Alph*)]
\item therefore \item in addition \item nonetheless \item because
\end{enumerate}
\end{enumerate}

\subsection*{Answer Key}

\begin{center}
\begin{tabular}{cll}
\toprule
\# & Ans & Explanation \\
\midrule
1 & B & ``Although'' + clause (has verb ``resembles'') \\
2 & C & Contrast between two sentences $\to$ ``however'' \\
3 & B & ``Despite'' + noun phrase (``its small market share'') \\
4 & B & ``So...that'' = result clause \\
5 & B & Contrast: supports BUT warns $\to$ ``nevertheless'' \\
6 & A & ``Despite AI IS'' $\to$ ``Although AI is'' or ``Despite AI being'' \\
7 & B & ``It's'' $\to$ ``its'' (possessive, not contraction) \\
8 & A & Inversion: ``Not only IS elementary OS beautiful'' \\
9 & A & ``Because'' + clause with verb (``evolves'') \\
10 & C & Contrast: values freedom BUT some prefer simplicity \\
\bottomrule
\end{tabular}
\end{center}

% ============================================================
\section*{Shadowing (Gemini Live --- 15 min)}
% ============================================================

\begin{quote}\small
``The question is not whether AI is intelligent. The question is
whether we will remain intelligent enough to use it wisely. When
we outsource our thinking to machines, we are not merely saving
time. We are atrophying the very cognitive muscles that make us
capable of critical thought.''
\end{quote}

\textbf{Targets:} \textbf{outsource} (OUT-source), \textbf{atrophying}
(AT-ro-fee-ing), \textbf{cognitive} (COG-ni-tiv), \textbf{merely}
(MEER-lee).

\end{document}